
By \SGARefLink{sga3:xvii:lemma:5:2:3:item:ii}{5.2.3(ii)}, \(R_n\) is representable by a nonempty affine
\(k\)-scheme.  Since a filtered direct limit of nonzero rings is
nonzero, the functor
\[
  R=\varprojlim_n R_n
\]
is representable by a nonempty affine \(k\)-scheme
(EGA IV, 8 and 1.9.1).  There are therefore an extension \(K\) of
\(k\) and a point \(u\in R(K)\).  The image \(u_n\) of \(u\) in
\(R_n(K)\) corresponds to a lift \(H'_n\) of \(({}_nH)_K\) in
\((E_n)_K\).  By construction,
\[
  H'_n={}_n(H'_m)\qquad\text{if \(n\) divides \(m\)}.
\]
Put
\[
  U_n=(U_K)^{H'_n}.
\]
The choice of \(H'_n\) identifies \((R_n)_K\) with \(U_K/U_n\).
The family of groups \(H'_n\) is filtered increasing, and hence
the family \(U_n\) is filtered decreasing.  It is therefore
stationary for sufficiently large \(n\), since \(U_K\) is
noetherian.  It follows that the family \((R_n)_K\) is stationary,
and consequently so is the family \(R_n\).  In short,
\[
  R_n=R
\]
for sufficiently large \(n\).

Under the hypotheses in force, \SGARefLink{sga3:xvii:lemma:5:2:3:item:i}{5.2.3(i)} shows that \(R_n(k)\) is
nonempty.  We may therefore find a coherent system of lifts
\(H'_n\) of \({}_nH\), for \(n\in N'\).  Let \(H'\) be the
smallest closed subscheme of \(E\) containing \(H'_n\) for every
\(n\) (Exposé VI\textsubscript{B}, \S7).  The argument of
\SGARefLink{sga3:xv:lemma:4:6}{Exposé XV, 4.6} shows that \(H'\) is a smooth commutative
algebraic group whose formation commutes with extension of the
base field.  To prove that \(H'\) is a lift of \(H\) in \(E\), we
may therefore suppose that \(k\) is algebraically closed.
According to BIBLE 4, Theorem 4, \(H'\) is then the direct product
of a smooth group \(M\) of multiplicative type and a unipotent
group \(V\).  The groups \(H'_n\) are necessarily contained in
\(M\) by \SGARefLink{sga3:xvii:proposition:2:4}{2.4}; the definition of \(H'\) then gives \(H'=M\).
Thus \(H'\) is of multiplicative type, and hence
\[
  H'\cap U=0.
\]
The morphism \(H'\to H\) is consequently a monomorphism, and the
density theorem of \SGARefLink{sga3:ix:remark:4:10}{Exposé IX, 4.10} shows that it is an
epimorphism.  It is therefore an isomorphism.

Now let \(H'\) and \(H''\) be two lifts of \(H\) in \(E\).  For
each \(n\in N'\), let
\[
  T_n=\operatorname{Transp}({}_nH',{}_nH'')
\]
be the transporter of \({}_nH'\) into \({}_nH''\).  It is
representable by a closed subscheme of \(U\)
(\SGARefLink{sga3:viii:examples:6:5:item:e:01}{Exposé VIII, 6.5(e)}).  The \(T_n\) form a filtered decreasing
family of closed subschemes of \(U\), and they are nonempty by
\SGARefLink{sga3:xvii:lemma:5:2:3:item:i}{5.2.3(i)}\SGARefLink{sga3:xvii:appendix:A}{(a)}.  Let \(T\) be their stationary value.  Under the
hypotheses of \SGARefLink{sga3:xvii:lemma:5:2:3:item:i}{5.2.3(i)}, \(T_n(k)\) is nonempty.  Thus there is an
element \(u\in U(k)\) such that
\[
  {}_nH''=\operatorname{int}(u)({}_nH')
  \qquad\text{for every \(n\in N'\)}.
\]
It follows that
\[
  H''=\operatorname{int}(u)H'
\]
by \SGARefLink{sga3:ix:corollary:4:8:item:b}{Exposé IX, 4.8(b)}.

\SGASubsection{5.4}{\texorpdfstring{The Case Where \(U\) Is
Radicial}{The Case Where U Is Radicial}}
\SGARefTarget{sga3:xvii:subsection:5:4}

\SGAStatement{Proposition 5.4.1}{If \(k\) is a perfect field of
characteristic \(p>0\), and if \(U\) is radicial, then the
extension \(E\) of \SGARefLink{sga3:xvii:theorem:5-1-1}{5.1.1} is trivial.}
\SGARefTarget{sga3:xvii:proposition:5:4:1}

Using a characteristic composition series of \(U\), we may
restrict ourselves to the case
\[
  U=(\boldsymbol\alpha_p)^r
\]
by \SGARefLink{sga3:xvii:lemma:3-9}{3.9}.

It follows from Appendix II, \SGARefLink{sga3:ii:subsection:2-2}{2.2} and \SGARefLink{sga3:xvii:proposition:2:1}{2.1} that there are
isomorphisms of \(k\)-functors
\[
  \operatorname{Aut}_{k\text{-}\mathrm{gr}}
    \bigl((\boldsymbol\alpha_p)^r\bigr)
  \xrightarrow{\ \sim\ }
  \operatorname{Aut}_{p\text{-}\mathrm{Lie}}
    \bigl(\operatorname{Lie}((\boldsymbol\alpha_p)^r)\bigr)
  \xrightarrow{\ \sim\ }
  \operatorname{GL}
    \bigl(\operatorname{Lie}((\boldsymbol\alpha_p)^r)\bigr).
\]
Let \(V\) be a \(k\)-vector space of rank \(r\), let
\(\mathbf W(V)\) be the corresponding vector \(k\)-group
(\SGARefLink{sga3:i:subsection:4-6}{Exposé I, 4.6}), and identify \((\boldsymbol\alpha_p)^r\) with
\({}_F V\).  The preceding remarks show that the action of \(H\)
on \({}_F V\) defined by the extension \(E\) comes from a linear
representation \(v\) of \(H\) on \(V\).  Consider the exact
sequence
\[
  \tag{*}\label{sga3:xvii:proposition:5-4-1:equation:star}
  0\longrightarrow {}_F V\longrightarrow V
    \xrightarrow{F}V^{(p)}\longrightarrow0,
\]
where \(F\) is the Frobenius morphism.
